\section{Элементы теории устойчивости.}
\subsection{Метод функций Ляпунова.}
Пусть дана некоторая система дифференциальных уравнений первого порядка в нормальной форме:
\begin{equation*}
    \dot{\bt{x}} = \bt{F}(\bt{x}, t).
\end{equation*}
Будем считать, что $\bt{F}$~---~гладкое векторное поле, определённой на всём $\R^n$.
\begin{definition}
    Будем называть точку $\bt{x}^0$ положением равновесия если система уравнений допускает постоянное решение $\bt{x}(t) = \bt{x}^0$.
\end{definition}
\begin{definition}
    Будем называть положение равновесия $\bt{x}^0 = 0$ устойчивым по Ляпунову если
    \[
        \forall \epsilon \exists \delta > 0: (\bt{x}(0), \dot{\bt{x}}(0)) \in U_\delta(0) \Ra \forall t > 0 (\bt{x}(t), \dot{\bt{x}}(t)) \in U_\epsilon(0).
    \]
    где под парой $(\bt{x}, \dot{\bt{x}})$ понимается некоторое решение системы с начальными условиями в $\delta$-окрестности нуля.
\end{definition}
\begin{note}
    Стоит заметить, что в определении устойчивого по Ляпунову положения равновесия явно предполагается бесконечная продолжаемость решений вправо.
\end{note}
\begin{note}
    Также заметим, что можно определить такое понятие как устойчивое решение системы дифференциальных уравнений -- аналогичным образом.
    Но если $\bt{x}^0(t)$~---~частное решение системы, то заменой $\bt{x} = \bt{y} + \bt{x}^0$ мы сводим задачу исследования решения системы на устойчивость к исследованию положения равновесия.
    В дальнейшем не будет делаться специальных оговорок и будут рассматриваться случаи положений равновесия.
\end{note}
\begin{note}
    Устойчивость положения равновесия не является инвариантном при замене координат.
    Пусть
    \[
        \dot{x} = \frac{1}{2}.
    \]
    Общим решением системы является $x(t) = \frac{t}{2} + x^0$.
    При малом возмущении решения
    \[
        \tilde{x}(t) = x^0 + \delta x^0 + \frac{t}{2}.
    \]
    видно, что
    \[
        |\tilde{x} - \tilde{x}| = \delta x^0.
    \]
    и решение системы устойчиво. \\
    Сделаем замену координат следующего вида:
    \[
        y = x^2.
    \]
    Решением системы будет
    \[
        y(t) = (x^0 + \frac{t}{2})^2.
    \]
    И его малым возмущением
    \[
        \tilde{y}(t) = (x^0 + \delta x^0 + \frac{t}{2})^2.
    \]
    Нетрудно заметить что
    \[
        |\tilde{y} - y| \geq \delta x^0 (x^0 + \frac{t}{2}).
    \]
    неограниченно растёт и, следовательно, в новых координатах $y(t)$ не является устойчивым решением.
\end{note}
Известно, что справедлива следующая теорема:
\begin{theorem}
    Пусть $\bt{x} = 0$ является устойчивым решением системы $\dot{\bt{x}} = \bt{f}(\bt{x}, t)$.
    Невырожденная замена координат $\bt{x} = \bt{F}(\bt{y}, t)$ такая, что как $\bt{F}(\bt{y}, t)$, так и $\bt{F}^{-1}(\bt{x}, t)$ равномерно непрерывны в нуле по $t$ и $\bt{F}(\bt{0}, t) = 0$ сохраняет устойчивость в том смысле, что $\bt{y} = 0$ является устойчивым решением системы
    \[
        \dot{\bt{y}} = D\bt{F}^{-1}\bt{f}(\bt{F}(\bt{y}, t), t).
    \]
\end{theorem}
\begin{proof}
    Заметим, что определение устойчивости по Ляпунову является определением равномерной непрерывности решений задачи Коши по времени в некоторой окрестности положения равновесия.
\end{proof}

